\chapter{Giữ Ngày Mai Từ Hôm Nay}
\markboth{Giữ Ngày Mai Từ Hôm Nay}{Protect Tomorrow Starting Today}

\section{Cuốn Sổ Cũ Trước Giờ Lên Lớp}
\begin{truyen}{Đời Người Đổi Hướng Bằng Thói Quen}{A Life Changes Direction Through Habits}
\chuhoa{S}{áng sớm, trước khi lên lớp, người thầy mở cuốn sổ cũ và đọc lại vài dòng mình đã ghi suốt nhiều tháng. Có dòng nhắc vì sao đã không mua một món đồ vô ích. Có dòng ghi số tiền nhỏ vừa dành vào quỹ dự phòng.}
\textit{(Early in the morning, before going to class, the teacher opened an old notebook and reread a few lines written over many months. Some lines reminded him why he had refused a useless purchase. Others recorded small sums recently added to the emergency fund.)}

Nhìn cuốn sổ, anh thấy một điều rất giản dị nhưng rất khó: đời người ít khi thay đổi vì một lời hứa thật lớn. Nó thường đổi hướng vì những việc nhỏ được lặp lại đủ lâu. Một lần ghi chép, một lần bớt tiêu, một lần từ chối sĩ diện, một lần đặt tiền sang bên cho ngày mai. Từng việc nhỏ không gây tiếng động, nhưng cộng lại thành một lối sống mới.
\textit{(Looking at the notebook, he saw something simple yet difficult: a life rarely changes because of one grand promise. It changes direction through small acts repeated for long enough. One act of recording, one act of spending less, one refusal of pride, one small amount set aside for tomorrow. None of these makes noise, but together they become a new way of living.)}

Điều làm anh ấm lòng nhất không phải con số lớn lên nhanh chóng. Điều ấy đến chậm. Cái quý hơn là cảm giác mình đã bắt đầu giữ được mình khỏi những cơn tiêu tiền vô định như trước.
\textit{(What warmed him most was not a rapidly growing number. That comes slowly. More precious was the feeling that he had begun to keep himself from the aimless spending storms that once controlled him.)}
\end{truyen}

\begin{baihoc}
Thay đổi tài chính bền vững thường đến từ các thói quen âm thầm, không phải từ các lời hứa ồn ào.
\textit{(Sustainable financial change usually comes from quiet habits, not loud promises.)}
\end{baihoc}

\begin{ghinhoanh}
Quiet habits change loud problems.

Repeat the right small things.
\end{ghinhoanh}

\begin{tuhocnhanh}
1. Thói quen nhỏ nào đang cứu em nhiều nhất? \textit{Which small habit is helping you most?}

2. Em có đang chờ một bước ngoặt lớn thay vì làm việc nhỏ mỗi ngày không? \textit{Are you waiting for a big turning point instead of doing small things daily?}
\end{tuhocnhanh}
\ngancach

\section{Thiên Lý Chi Hành}
\begin{truyen}{Một Bước Chân Cũng Là Bắt Đầu}{Even One Step Is a Beginning}
\chuhoa{L}{ão Tử nói: ``Thiên lý chi hành, thủy ư túc hạ'' --- hành trình ngàn dặm bắt đầu từ dưới chân mình. Câu này ai cũng từng nghe, nhưng ít ai áp vào chuyện tiết kiệm với sự kiên nhẫn thật sự.}
\textit{(Laozi said, ``A journey of a thousand miles begins beneath one’s feet.'' Everyone has heard this line, but few apply it to saving with real patience.)}

Người ta thường đợi khi nào có nhiều tiền mới bắt đầu để dành, khi nào bớt khó khăn mới bắt đầu kỷ luật, khi nào đỡ bận mới bắt đầu ghi chép. Nhưng đời ít khi tự sắp cho mọi thứ thuận rồi mới cho ta bắt đầu. Chính trong hoàn cảnh còn bận, còn chật, còn thiếu, một bước nhỏ mới có giá trị lớn hơn hết.
\textit{(People often wait until they have more money to start saving, until hardship lessens to begin discipline, until busyness fades to begin recording expenses. But life rarely arranges everything conveniently before allowing us to begin. It is precisely in the midst of busyness, narrowness, and lack that a small step has the greatest value.)}

Người thầy hiểu rằng nếu cứ đợi đủ điều kiện, anh sẽ mãi đứng yên. Vì thế, anh chọn đi bằng bước nhỏ: một khoản để dành ít thôi, một câu tự nhắc ngắn thôi, một cuốn sổ đơn giản thôi. Bước nhỏ nhưng có thật tốt hơn mọi kế hoạch đẹp mà không làm nổi.
\textit{(The teacher understood that if he kept waiting for ideal conditions, he would stand still forever. So he chose to move with small steps: a modest amount saved, a short self-reminder, a simple notebook. A small real step is better than every beautiful plan that is never carried out.)}
\end{truyen}

\begin{baihoc}
Bắt đầu nhỏ không đáng sợ. Đáng sợ là cứ đợi mãi rồi không bắt đầu.
\textit{(Beginning small is not frightening. What is frightening is waiting forever and never beginning.)}
\end{baihoc}

\begin{ghinhoanh}
Start small.

Reality is better than postponed perfection.
\end{ghinhoanh}

\begin{tuhocnhanh}
1. Việc nhỏ nào em có thể bắt đầu ngay hôm nay? \textit{What small action can you begin today?}

2. Em đang đợi điều kiện nào rồi mới chịu thay đổi? \textit{What condition are you waiting for before allowing yourself to change?}
\end{tuhocnhanh}
\ngancach

\section{Mỗi Lần Bớt Một Khoản Là Một Lần Giữ Lấy Tương Lai}
\begin{truyen}{Không Có Chiến Công Lớn, Chỉ Có Sự Bền Bỉ}{No Grand Heroics, Only Steadiness}
\chuhoa{N}{gười thầy khép cuốn sổ lại, uống một ngụm trà, rồi đi dạy. Trong lòng anh không phải không còn lo, nhưng nỗi lo ấy đã bớt mù mờ vì mỗi ngày đều có một việc nhỏ để chống lại nó.}
\textit{(The teacher closed the notebook, drank a sip of tea, and went to teach. His heart was not free from worry, but the worry had grown less shapeless because each day now contained one small act resisting it.)}

Anh hiểu rằng mình không cần một chiến công tài chính oanh liệt. Anh cần sự bền bỉ: bớt một khoản lãng phí, giữ một khoản dự phòng, từ chối một lần mua sắm vì sĩ diện, dạy con thêm một bài học nhỏ về biết đủ. Từng điều ấy chính là cách giữ lấy ngày mai.
\textit{(He understood that he did not need some heroic financial triumph. He needed steadiness: reducing one wasteful expense, protecting one emergency fund, refusing one pride-driven purchase, teaching his child one more small lesson about enoughness. Each of these was a way of protecting tomorrow.)}

Ngày mai bình yên không được gửi từ trên trời xuống. Nó được góp thành từ những ngày hôm nay biết chắt chiu. Và người biết chắt chiu không phải là người sống hẹp, mà là người đã hiểu cái giá của những lần buông tay.
\textit{(A peaceful tomorrow is not delivered from the sky. It is assembled from today’s acts of careful preservation. And the one who preserves carefully is not narrow-hearted, but someone who has learned the price of letting go.)}
\end{truyen}

\begin{baihoc}
Mỗi lần tránh một khoản lãng phí là một viên gạch nhỏ đặt cho tương lai bớt chông chênh.
\textit{(Each avoided waste is a small brick laid for a steadier future.)}
\end{baihoc}

\begin{ghinhoanh}
Protect tomorrow today.

Steadiness builds shelter.
\end{ghinhoanh}

\begin{tuhocnhanh}
1. Ngày mai của em đang được xây bởi thói quen nào hôm nay? \textit{Which habit today is building your tomorrow?}

2. Viên gạch nhỏ đầu tiên em muốn đặt tiếp là gì? \textit{What is the next small brick you want to lay?}
\end{tuhocnhanh}

\begin{turanminh}
1. Tôi không cần chờ một bước ngoặt lớn mới sửa đời mình; tôi chỉ cần bền bỉ với những việc nhỏ đúng đắn mỗi ngày.

2. Mỗi khoản lãng phí tránh được, mỗi dòng ghi chép giữ lại, mỗi lần dừng tay đúng lúc đều là một viên gạch cho ngày mai của gia đình.

3. Tôi đã đi lên từ chỗ rất thấp, nên càng phải sống sao cho con mình không phải đi lại con đường quá gập ghềnh ấy.
\end{turanminh}

\begin{dayhoc}
1. Chương cuối có thể dùng để nhắc học sinh rằng tương lai không được xây bằng hứng khởi nhất thời, mà bằng nếp sống lặp lại mỗi ngày.

2. Thầy có thể yêu cầu các em viết một ``viên gạch nhỏ'' cho tuần tới: một việc đơn giản nhưng đủ thật để làm đời mình tốt hơn.

3. Câu kết phù hợp trên lớp là: ``Giữ ngày mai không phải chuyện lớn lao; đó là biết sống đúng từ hôm nay.''
\end{dayhoc}